\documentclass[11pt]{article}
\usepackage[margin=25mm]{geometry}
\usepackage{amsmath,amssymb,amsthm,mathtools,graphicx,hyperref,xurl}
\hypersetup{hypertexnames=false,colorlinks=true,urlcolor=blue}
\newtheorem{theorem}{Theorem}
\newtheorem{proposition}[theorem]{Proposition}
\newcommand{\C}{\mathbb C}
\newcommand{\dd}{\mathfrak D_{\rm odd}}
\newcommand{\rem}{\operatorname{rem}}
\newcommand{\diag}{\operatorname{diag}}
\title{A receiving-frame divisor and its exact one-moment repair}
\author{Split-Zero programme: independently checked continuation of received web-session work}
\date{20 September 2026}
\begin{document}\maketitle

\section{The original maps and the scope of this calculation}
The original four-coordinate Fable receiving map has eight marked states
over a squarefree quartic. This note proves that its particular odd
evaluation frame can become singular while all eight states remain
regular, and calculates a minimal observation that recovers the lost
direction. The original source's arithmetic-quartet invertibility
remains valid: its exact pullback is proved below, with the original
leading coefficient retained. The general target used to exhibit the
new divisor is not asserted to be a quartet of zeta zeros.

The received formulas and proposed repair are attributed to the
user-supplied web-session calculation \cite{Incoming}; the original
marked-state construction is \cite[SE1--14, GF20--21]{Fable}%
% BEGIN VERIFIED PUBLIC PROOF LINK
~\href{https://github.com/KokunoYumeto/zeta-function-research-reader/blob/5992ad50c0f2a06921f1fcaded7b4e38097c0739/workbenches/splitzero-tandem/continuations/20260920-original-kernel-mixed-determinants/providers/FABLE_TO_ORIGINAL_CONDUCTOR.tex\#L1598}{[proof GF20]}%
% END VERIFIED PUBLIC PROOF LINK
{}.
The classical interpolation used in the null-vector calculation is
Lagrange interpolation, in the original equation sources of
N.~M.~Temme \cite{Temme}. Every polynomial identity needed here is
also proved below; the accompanying exact rational script verifies
their coefficient equalities without substituting numerical roots.

Let $u=(A,B,C,D)\in\C^4$, with $A\ne0$, and
\begin{equation}
 h(r)=Ar^4+r^3+Br^2+Cr+D,\quad d(r)=h'(r).
 \tag{OF1}
\end{equation}
Assume $h$ has four distinct roots in a fixed order $r_1,\ldots,r_4$.
Choose and retain $\xi_j^2d(r_j)=1$. The eight original columns are
$p_j^\pm=e_j\pm o_j$, where
\begin{equation}
 e_j=\begin{pmatrix}0\\-r_j\\3i d_j\\
 (B-17r_j+Ar_j^2)d_j-2Ad_j^2\end{pmatrix},\qquad
 o_j=\xi_j\begin{pmatrix}1\\-id_j\\Ad_j^2+2r_jd_j\\
 i(7r_j^2d_j-13d_j^2)\end{pmatrix},\quad d_j=d(r_j).
 \tag{OF2}
\end{equation}
For $E=[e_j]$, $O=[o_j]$, and literal branch coefficients $c_j^\pm$,
\begin{equation}
 \alpha_j=c_j^++c_j^-,\quad\beta_j=c_j^+-c_j^-,\quad
 \sum_j(c_j^+p_j^++c_j^-p_j^-)=E\alpha+O\beta,\quad
 c_j^\pm=(\alpha_j\pm\beta_j)/2.\tag{OF3}
\end{equation}
These formulas specify all domains and inverses between the marked
eight-dimensional coefficient space and its even/odd coordinates.
An original invertible linear conductor $L:\C^4\to\C^4$ composes
as $LO$ and satisfies $\det(LO)=\det L\det O$. Thus the receiving
divisor computed here is carried through that conductor unchanged.

\section{The complete divisor and its arithmetic pullback}
\begin{theorem}
Put
\begin{align}
 \dd={}&20(3C-B^2)+A(87B^3-282BC-51D)\notag\\
 &+A^2(-28B^4+22B^2C+431C^2+204BD)\notag\\
 &+A^3(224B^2D-168BC^2-856CD)-448A^4D^2.
 \tag{OF4}
\end{align}
Then, for the retained branch ordering and signs,
\begin{equation}
 \det O=\varepsilon_{\rm br}\frac{4\dd}{A^3},\qquad
 \varepsilon_{\rm br}\in\{1,-1\}.\tag{OF5}
\end{equation}
\end{theorem}
\begin{proof}
Reduce the four polynomials in $o_j/\xi_j$ modulo $h$, retaining
the power basis $(1,r,r^2,r^3)$. Their coefficient matrix is
\begin{equation}
 \begin{aligned}
 Q_{1,*}&=(1,0,0,0),\\
 Q_{2,*}&=(-iC,-2iB,-3i,-4iA),\\
 Q_{3,*}&=(AC^2-9D,\ 4ABC-8AD-7C,\\
 &\hspace{14mm}4AB^2-2AC-16A^2D-5B,\ 4AB-8A^2C-3),\\
 Q_{4,*}&=i(-13C^2+20D/A,\ 20C/A+76D-52BC,\\
 &\hspace{14mm}20B/A+5C-52B^2+208AD,\ 20/A-66B+104AC).
 \end{aligned}\tag{OF6}
\end{equation}
For example $Ar^4=-r^3-Br^2-Cr-D$ reduces every higher power;
successive substitutions give precisely the four rows displayed.
Expanding the remaining three-by-three determinant after the first
row gives $\det Q=4\dd/A$. If
$V=(r_j^{k})_{0\le k\le3,1\le j\le4}$ and
$\Delta=\det V=\prod_{j<k}(r_k-r_j)$, then
$O=QV\diag(\xi_j)$. Also
\[
 \prod_jd_j=A^4\prod_j\prod_{k\ne j}(r_j-r_k)
 =A^4\Delta^2.
\]
There are six reversed pairs, giving positive sign. Therefore
$(\Delta\prod_j\xi_j)^2=A^{-4}$, which proves OF5.
\end{proof}

On the original arithmetic quartet, with the parameters of
\cite[GF20]{Fable}%
% BEGIN VERIFIED PUBLIC PROOF LINK
~\href{https://github.com/KokunoYumeto/zeta-function-research-reader/blob/5992ad50c0f2a06921f1fcaded7b4e38097c0739/workbenches/splitzero-tandem/continuations/20260920-original-kernel-mixed-determinants/providers/FABLE_TO_ORIGINAL_CONDUCTOR.tex\#L1598}{[proof GF20]}%
% END VERIFIED PUBLIC PROOF LINK
{}, substitution retains
\begin{equation}
 (A,B,C,D)=\left(-\tfrac12,\delta^2-\gamma^2-\tfrac34,
 \gamma^2-\delta^2+\tfrac14,
 -\tfrac1{32}-\tfrac{\gamma^2-\delta^2}{4}
 -\tfrac{(\delta^2+\gamma^2)^2}{2}\right).
 \tag{OF7}
\end{equation}
Expansion of OF4 gives
\[
 \dd(u_\rho)=-56\delta^2\gamma^2
                 (2\delta^2\gamma^2-\delta^2+\gamma^2).
\]
Thus $\det Q=448\delta^2\gamma^2
(2\delta^2\gamma^2-\delta^2+\gamma^2)>0$ for the source domain
$\delta\ne0$, $\gamma>2$. This proves the precise relation to the
existing quartet theorem. No new localization is necessary there.

\section{A regular target where one receiving direction is lost}
Set $A=1$, $B=D=0$ and retain $C=c$. Then
$h_c(r)=r(r^3+r^2+c)$ and $\dd=c(431c+60)$.
Write $c_*=-60/431$ and $h_*=h_{c_*}$. This quartic is squarefree:
its root zero is simple, and a repeated root of its cubic factor
would have to be $0$ or $-2/3$. The cubic at those points is $c_*$
or $4/27+c_*=104/11637$, respectively, both nonzero.

Let $p_0,p_1,p_2,p_3$ be the original component polynomials in OF2.
For any polynomial $F$, Lagrange interpolation gives the exact
residue identity
\begin{equation}
 \sum_j\frac{F(r_j)}{h_*'(r_j)}=[r^3]\rem_{h_*}F.\tag{OF8}
\end{equation}
Indeed the remainder $R$ has degree at most three and equals
$\sum_jR(r_j)h_*(r)/((r-r_j)h_*'(r_j))$ by evaluation at its four
distinct nodes. The leading coefficient of each quotient is one.
This is the specialization of \cite[3.3.1--3.3.3\_1]{Temme} used here.

Define
\begin{equation}
 \ell^T=(208c_*^2,5ic_*,60,-21i),\quad
 q_*(r)=973r^2+413r-140,\quad b_j=\xi_jq_*(r_j).
 \tag{OF9}
\end{equation}
Direct reduction by $h_*$ gives
\[
 \rem_{h_*}\Big(\sum_{j=0}^3\ell_jp_j\Big)=0,
 \qquad [r^3]\rem_{h_*}(q_*p_j)=0\quad(0\le j\le3).
\]
The first identity proves $\ell^TO=0$, and OF8 proves $Ob=0$.
The vector $b$ is nonzero because a nonzero polynomial of degree two
cannot vanish at four distinct nodes and all $\xi_j\ne0$.

For an exact rank certificate introduce the frame
\begin{equation}
 \widetilde O_{\cdot j}=\xi_j(1,-id_j,d_j^2+2r_jd_j,r_j^3)^T,
 \quad
 L_c=\begin{pmatrix}
 1&0&0&0\\0&1&0&0\\0&0&1&0\\
 -208ic^2/21&5c/21&-20i/7&4i(431c+60)/21
 \end{pmatrix}.\tag{OF10}
\end{equation}
Reduction of the last original component proves, coefficient by
coefficient,
\begin{equation}
 O=L_c\widetilde O,\qquad
 \det\widetilde O=\varepsilon_{\rm br}(-21ic),\qquad
 \widetilde O(c_*)b=420e_4.\tag{OF11}
\end{equation}
For clarity the last component identity is
\[
 p_3\equiv-\frac{208ic^2}{21}p_0+\frac{5c}{21}p_1
 -\frac{20i}{7}p_2+\frac{4i(431c+60)}{21}r^3\pmod{h_c}.
\]
The determinant follows by expanding the coefficient matrix of the
four polynomials defining $\widetilde O$: it equals $-21ic$, and
the retained root/sign factor is $\varepsilon_{\rm br}$ since $A=1$.
For the final equality in OF11, OF8 applied to
$q_*,q_*p_1,q_*p_2,q_*r^3$ gives coefficients $0,0,0,420$.
The first three rows of $L_{c_*}$ are independent and its last
diagonal entry is zero. Thus $O(c_*)$ has rank exactly three.

\section{The pole, its metric coefficient and a finite remainder}
Analytic roots and square roots exist on a sufficiently small disk
about $c_*$, since the discriminant has no zero there. Retain the
initial marking and its analytic continuation. Put $\epsilon=c-c_*$.
Block inversion of OF10 gives
\begin{equation}
 \lim_{c\to c_*}\epsilon L_c^{-1}=e_4\ell^T/1724,
 \qquad
 O(c)^{-1}=\frac{b\ell^T}{724080\epsilon}+O(1).
 \tag{OF12}
\end{equation}
The factor $724080=420\cdot1724$ follows from
$\widetilde O(c_*)^{-1}e_4=b/420$. This proves the full residue,
including all phases. For fixed positive Hermitian matrices $R,G$,
\begin{equation}
 \|O(c)^{-1}\|_{G\to R}
 =\frac{\sqrt{b^*Rb}\sqrt{\ell^TG^{-1}\bar\ell}}
 {724080|\epsilon|}+O(1).\tag{OF13}
\end{equation}
To verify this, conjugate the rank-one residue by $R^{1/2}$ and
$G^{-1/2}$. Its Euclidean operator norm is the product of the
two vector norms. The bounded remainder changes the norm by a
bounded amount by the triangle inequality in both directions.
For a fixed invertible original conductor $L$ this also gives
\[
 (LO(c))^{-1}=\frac{b\ell^TL^{-1}}{724080\epsilon}+O(1).
\]

The remainder may be bounded explicitly. Define
\begin{gather}
 R_0=52/11637,\quad r_0=R_0/2,\quad m_c=1568/11637,
 \quad\delta_0=m_c^3(52/431),\notag\\
 M=\max\{512/\sqrt{\delta_0},6880\},\quad B_{
 \rm adj}=24M^3,\quad d_{\rm det}=6376/27,\notag\\
 C_{\rm rem}=\frac{B_{\rm adj}}{d_{\rm det}}
                   \left(\frac2{R_0}+\frac{1724}{240}\right).
 \tag{OF14}
\end{gather}
Then
\begin{equation}
 \left\|O(c)^{-1}-\frac{b\ell^T}{724080\epsilon}\right\|_2
 \le C_{\rm rem},\qquad 0<|\epsilon|\le r_0.\tag{OF15}
\end{equation}
Here are the bounds retaining their constants. On $|\epsilon|\le R_0$,
$|c|\ge m_c$, $|4+27c|\ge52/431$, and $|c|<1$.
The nonzero roots solve $r^3+r^2+c=0$; if $|r|\ge2$ then
$|r|^3>|r|^2+|c|$, a contradiction. Thus all roots have modulus
less than two and $|d_j|\le64$. Also
$|\prod d_j|=|c|^3|4+27c|\ge\delta_0$, so
$|d_j|\ge\delta_0/64^3$. The first row of $O$ is bounded by
$512/\sqrt{\delta_0}$. The next three rows have respective bounds
$8$, $64^{3/2}+4\sqrt{64}=544$, and
$28\sqrt{64}+13\cdot64^{3/2}=6880$. Each adjugate entry has six
terms bounded by $M^3$, hence its operator norm is at most $24M^3$.
All branches extend analytically over the closed disk and slightly
beyond it. Cauchy's power-series coefficient bound therefore gives
\[
 \|\operatorname{adj}O(c)-\operatorname{adj}O(c_*)\|_2
 \le 2B_{\rm adj}|\epsilon|/R_0\quad(|\epsilon|\le R_0/2).
\]
Finally $\det O=\varepsilon_{\rm br}\epsilon(-240+1724\epsilon)$,
and $|-240+1724\epsilon|\ge d_{\rm det}$ on this smaller disk.
Subtracting the constant adjugate residue in this exact inverse
formula proves OF15. For the $G\to R$ norm multiply the right side
by $\|R^{1/2}\|_2\|G^{-1/2}\|_2$.

\section{One retained scalar repairs the complete marked observation}
Let $m=(\xi_jr_j^3)_j$ and choose an explicit weight $\eta>0$.
For the original $G>0$ the augmented odd observation has Gram
\begin{equation}
 \beta\longmapsto (O\beta,m\beta),\qquad
 H_\eta=O^*GO+\eta m^*m.\tag{OF16}
\end{equation}
It is injective near and at $c_*$: retain the first three coordinates
of $O\beta$ and append $m\beta$, obtaining $\widetilde O\beta$.
Invert OF11. The original fourth coordinate is then recovered by
$L_c$. This single additional scalar is minimal: the old kernel
has dimension one, and $m(c_*)b=420$.

Write $K_c$ for the first three columns of $L_c$. Then
\begin{equation}
 \det H_\eta=16|c|^2|431c+60|^2\det G
       +441\eta|c|^2\det(K_c^*GK_c).\tag{OF17}
\end{equation}
Indeed OF11 turns the Gram into
$\widetilde O^*(L_c^*GL_c+\eta e_4e_4^*)\widetilde O$.
The determinant is affine in the added $(4,4)$ entry; its
coefficient is the top three-by-three minor $\det(K_c^*GK_c)$.
Now $|\det\widetilde O|^2=441|c|^2$ and
$|\det L_c|^2=16|431c+60|^2/441$. This proves OF17.
At $c_*$ it is strictly positive because the three columns of $K_c$
are independent and $G>0$. No original Gram has been replaced.

The full original eight-coordinate map is
\begin{equation}
 J(\alpha,\beta)=(\alpha,E\alpha+O\beta).
 \tag{OF18}
\end{equation}
At $c_*$ it has kernel $\{(0,tb):t\in\C\}$ and rank seven.
Let $P_3:\C^4\to\C^3$ select its first three entries and set
$\widetilde E=(P_3E;0)$. Its repaired eight-coordinate map is
\begin{equation}
 \widetilde J(\alpha,\beta)
 =\left(\alpha,\begin{pmatrix}P_3(E\alpha+O\beta)\\m\beta\end{pmatrix}\right)
 =\begin{pmatrix}I&0\\\widetilde E&\widetilde O\end{pmatrix}
   \binom\alpha\beta.\tag{OF19}
\end{equation}
For output $(\alpha,z)$ its inverse is
$\beta=\widetilde O^{-1}(z-\widetilde E\alpha)$, and the entire
old evaluation is $E\alpha+L_c(z-\widetilde E\alpha)$.
Thus the fourth original coordinate is recovered after subtracting
and restoring its full even contribution; that contribution has
not been dropped.

The actual programme target direct sum retains the two matrices
$G_Y$ and $G$ \cite[RW14]{Fable}%
% BEGIN VERIFIED PUBLIC PROOF LINK
~\href{https://github.com/KokunoYumeto/zeta-function-research-reader/blob/5992ad50c0f2a06921f1fcaded7b4e38097c0739/workbenches/splitzero-tandem/continuations/20260920-original-kernel-mixed-determinants/providers/FABLE_TO_ORIGINAL_CONDUCTOR.tex\#L1598}{[proof GF20]}%
% END VERIFIED PUBLIC PROOF LINK
{}. Its augmented metric is exactly
\begin{equation}
 \widehat H_\eta=
 \begin{pmatrix}I&0\\E&O\end{pmatrix}^{\!*}
 \diag(G_Y,G)
 \begin{pmatrix}I&0\\E&O\end{pmatrix}
 +\eta\begin{pmatrix}0&0\\0&m^*m\end{pmatrix}.
 \tag{OF20}
\end{equation}
It is positive definite throughout the repaired neighborhood:
zero norm first forces $\alpha=0$, then $O\beta=0$ and
$m\beta=0$, which force $\beta=0$ by OF16. This is the exact
return to the original metric, rather than an identification of
$G_Y$ with $G$. The source return uses its own retained pair
$G_X,G_U$ from \cite[RW19]{Fable}%
% BEGIN VERIFIED PUBLIC PROOF LINK
~\href{https://github.com/KokunoYumeto/zeta-function-research-reader/blob/5992ad50c0f2a06921f1fcaded7b4e38097c0739/workbenches/splitzero-tandem/continuations/20260920-original-kernel-mixed-determinants/providers/FABLE_TO_ORIGINAL_CONDUCTOR.tex\#L1598}{[proof GF20]}%
% END VERIFIED PUBLIC PROOF LINK
{} by the same formula.

\section{A finite inverse that extends after removal of its residue}
The patch also evaluates the formerly unspecified constant term.
Write $T(c)=\widetilde O(c)^{-1}$, $t(c)=T(c)e_4$ and
$v(c)=(-L_{41}(c),-L_{42}(c),-L_{43}(c),1)$.
With $a_0=1724i/21$ and $D_3=\diag(1,1,1,0)$, the exact identity is
\begin{equation}
 O(c)^{-1}=T(c)D_3+\frac{t(c)v(c)}{a_0(c-c_*)}.\tag{OF21}
\end{equation}
Consequently its constant Laurent coefficient is
\begin{equation}
 B_0=T_*D_3+
 \frac{-T_*\widetilde O'_*T_*e_4v_*+T_*e_4v'_*}{a_0},
 \quad
 v'_*=\left(\frac{416ic_*}{21},-\frac5{21},0,0\right).
 \tag{OF22}
\end{equation}
All entries are now finite sums on the actual four marked roots.
Specifically, $r_j'=-r_j/d_j$ and
$d_j'=1-(12r_j^2+6r_j)r_j/d_j$, whence
$\xi_j'=-\xi_jd_j'/(2d_j)$. Differentiating each entry of OF10
with these formulas evaluates $\widetilde O'_*$ without a new
limit or unspecified inverse. The identity $T'=-T\widetilde O'T$
follows by differentiating $T\widetilde O=I$ and proves OF22.
Thus the finite remainder is both bounded by OF15 and explicitly
evaluated at its limiting point, in the same original coordinates.

\clearpage
\begin{figure}[p]
\centering
\includegraphics[width=\linewidth]{ODD_FRAME_REPAIR_ILLUSTRATION.png}
\caption{An exact section of the receiving map and its repair. For
$\beta=\widetilde O(c)^{-1}(s,0,0,t)^T$, with real $s,t\in[-1,1]$,
the drawn coordinates are
$(y_1,\operatorname{Im}y_4,m\beta)
=(s,-208c^2s/21+4(431c+60)t/21,t)$.
Both targets have four simple roots. The dark coordinate projection
forgets the added cubic moment. It has rank two on the left, but
rank one at $c_*=-60/431$ on the right. The gold line is
$\beta=tb/420$; its original output is zero and its added moment
is $t$. OF9--20 prove the exact inverse, minimal repair and original
two-Gram return. This is a coordinate section, not an isometric
embedding of the Gamma norm. Receiving columns: \cite[SE1--14]{Fable}%
% BEGIN VERIFIED PUBLIC PROOF LINK
~\href{https://github.com/KokunoYumeto/zeta-function-research-reader/blob/5992ad50c0f2a06921f1fcaded7b4e38097c0739/workbenches/splitzero-tandem/continuations/20260920-original-kernel-mixed-determinants/providers/FABLE_TO_ORIGINAL_CONDUCTOR.tex\#L1598}{[proof GF20]}%
% END VERIFIED PUBLIC PROOF LINK
{};
classical null-vector interpolation: \cite[3.3(i)]{Temme}; received
divisor and repair: \cite{Incoming}.}
\end{figure}
\clearpage
\begin{thebibliography}{9}
\bibitem{Fable} Split-Zero programme, \emph{Fable to the original
conductor}, source \href{https://github.com/KokunoYumeto/zeta-function-research-reader/blob/5992ad50c0f2a06921f1fcaded7b4e38097c0739/workbenches/splitzero-tandem/continuations/20260920-original-kernel-mixed-determinants/providers/FABLE_TO_ORIGINAL_CONDUCTOR.tex\#L1}{\texttt{FABLE\_TO\_ORIGINAL\_CONDUCTOR.tex}},
consulted source hash
\nolinkurl{9629ec8ff32c069a7329c012687fd8bcfd7d5c9437a94fcfba1d4725edbec764}.
GF20--21, SE1--14 and RW14--19 are the exact receiving passages used.
\url{https://github.com/KokunoYumeto/zeta-function-research-reader}.
This programme source records the underlying human-source provenance
of the original polynomial separately; the receiving divisor and
repair here are attributed to the received continuation below.

% BEGIN VERIFIED PUBLIC PROOF LINK
\href{https://github.com/KokunoYumeto/zeta-function-research-reader/blob/5992ad50c0f2a06921f1fcaded7b4e38097c0739/workbenches/splitzero-tandem/continuations/20260920-original-kernel-mixed-determinants/providers/FABLE_TO_ORIGINAL_CONDUCTOR.tex\#L1598}{Source proof GF20}.
% END VERIFIED PUBLIC PROOF LINK
\bibitem{Incoming} User-supplied web-session calculation,
\emph{The relevant strand is Erd\H{o}s--Straus/Fable},
received 20 September 2026; no human author name supplied.
Original SHA256
\nolinkurl{57efa66b73bb4ce951b1d17ec4170b0f7bf1a00af3718e56faa71e7693ec2717}.
Equations (1)--(28) supply the divisor, pole and proposed repair;
OF19--22 articulate its full even-coordinate return and evaluate
the finite inverse coefficient. Original input remains preserved
in the private source register.
\bibitem{Temme} N.~M.~Temme, \emph{Numerical Methods}, NIST Digital
Library of Mathematical Functions, version 1.2.8, section 3.3(i),
equations 3.3.1, 3.3.2, 3.3.3 and 3.3.3\_1,
\url{https://dlmf.nist.gov/3.3.i}.
Original equation TeX read and retained. The DLMF source notes cite
P.~J.~Davis, \emph{Interpolation and Approximation} (1975),
pages 56 and 67--68; this note consulted the original DLMF equations,
not that entire book. OF8 gives the full proof of the specialization.
\end{thebibliography}
\end{document}
